\chapter*{Conclusion}
\addcontentsline{toc}{chapter}{Conclusion}
\markboth{Conclusion}{Conclusion}

This thesis set out from an ordinary observation: when an object is lost, the
information needed to recover it usually already exists, split between two
people who have no way of finding each other. The channels available today ---
unconnected lost-property desks, social-media groups where a report survives a
few hours, word of mouth --- fail for one structural reason rather than through
any lack of goodwill. Nothing in them is responsible for comparing a loss report
against a discovery report.

\projectname{} assigns that responsibility to software. We designed and built a
platform on which a single guided report enters a national pool that is searched
automatically and continuously, in both directions. The comparison is multimodal
and multilingual: descriptions are encoded by a multilingual sentence model, and
photographs by a vision-language model searched as an independent channel. Both
vector spaces live inside PostgreSQL through the \code{pgvector} extension, so
that nearest-neighbour retrieval and relational filtering are one query. Around
the engine, the prototype provides a trilingual interface with right-to-left
support, a claim workflow that protects contact details, a paid matching tier
built on an auditable credit ledger, and an administration console whose every
action is recorded.

The prototype was measured rather than only described. In a controlled
experiment, the multilingual encoder matched the right object first in all 72
cross-language queries, where an English-only model failed almost every query
involving Arabic. On the platform's own photographs, the raw similarity of two
unrelated images was found to have a median of $0.58$ --- enough, fed directly
into a weighted average, to put unrelated objects in front of users while every
unit test passed --- and the calibration of the image channel was set against
that measured distribution. Of the eight suggestions judged by their owners,
seven were correct; three of the confirmed ones were carried by their
photographs despite weak textual agreement, which is the case the independent
image route was built for. The whole platform runs in about 1.2\,GB of memory,
without a GPU, and answers its users in tens of milliseconds.

Those measurements also gave the business plan its numbers. The pilot costs
about 4\,900 dinars a month, more than half of it the levies of the legal status
rather than the servers. Selling match unlocks at 200 to 1\,200 dinars covers
those costs from a few hundred reports a month, but paying a team requires a
national volume. The project's remaining risk is therefore commercial ---
volume and conversion --- and not technical, and both are what a pilot will
measure.

Accuracy alone would not have produced a usable product. A high-confidence match
is a hypothesis about two objects, not an authorisation to disclose a person's
telephone number, and the platform separates the two completely: recovery
passes through verification questions set by the reporter, answered by the
claimant, and judged by the reporter. The one suggestion owners rejected in our
data had a perfect similarity score, which is the clearest argument for keeping
that human step.

\section*{Future work}
\addcontentsline{toc}{section}{Future work}

The system was built to accommodate the following directions without structural
change; they are the same as the future vision of Chapter~\ref{ch:innovation}.

\begin{itemize}[leftmargin=1.4em,itemsep=4pt]
  \item \textbf{A field pilot in Bouira.} Deploying with the university and a
        first partner institution, to measure what no demonstration corpus can:
        recovery rates, the share of owners who unlock a suggestion, and the
        share who buy.

  \item \textbf{Learned scoring.} Every owner verdict is stored next to the
        scores that produced the suggestion. That dataset is what is needed to
        replace the hand-set weights and boosts with a calibrated model, learned
        globally and per category --- photographs should dominate for bags,
        text for documents and keys.

  \item \textbf{Cross-modal matching.} Comparing the words of one report against
        the photograph of another, which the vision-language model already
        supports and which would rescue pairs where one side has only text and
        the other only an image.

  \item \textbf{Complete localisation.} Generating notifications as codes
        translated on the reader's side, as match explanations already are, and
        adding the eleven wilayas created in November 2025.

  \item \textbf{Institutional integration.} Letting partner lost-property desks
        publish their registers directly into the shared pool, possibly as the
        paid institutional accounts left open in Chapter~\ref{ch:business}.

  \item \textbf{Native mobile application}, enabling a report to be filed from
        the camera at the moment of discovery.
\end{itemize}
